
\chapter{Những Cuộc Trốn Chạy Bình Thường}
\markboth{Những Cuộc Trốn Chạy Bình Thường}{Những Cuộc Trốn Chạy Bình Thường}

\section{Bỏ Học Không Phải Lúc Nào Cũng Lười}
\begin{truyen}{Bỏ Học Không Phải Lúc Nào Cũng Lười}{Skipping Class Is Not Always Laziness}
\chuhoa{H}{ọc sinh bỏ tiết Lý liên tục ba tuần.} Không phải vì lười. Vì bạn thân ngồi cùng lớp vừa xảy ra chuyện nặng nề, và mỗi khi nhìn thấy nhau là không khí nghẹt lại. Ngồi trong lớp đó không học được gì.
\textit{(A student skips Physics three weeks in a row. Not out of laziness. Because a close friend in the same class just went through something serious, and whenever they see each other the air becomes unbearable. Sitting in that room makes learning impossible.)}

\teacherpair{Bỏ học liên tục là thiếu ý thức. Sẽ liên hệ phụ huynh.}{``Repeatedly skipping class shows a lack of responsibility. I will contact your parents.''}

Phụ huynh được gọi vào. Câu chuyện thật không bao giờ được kể. Vì không ai hỏi.
\textit{(The parents are called in. The real story is never told. Because nobody asked.)}

\begin{baihoc}
Hành vi là triệu chứng. Bỏ học, không nộp bài, ngủ gật trong lớp --- tất cả đều có thể là dấu hiệu của điều khác, lớn hơn nhiều. Hỏi trước khi phán xét không phải mềm lòng, đó là cách tiếp cận có thông tin hơn.
\textit{(Behaviour is a symptom. Skipping class, missing assignments, sleeping in class, all can be signs of something else, something much larger. Asking before judging is not weakness. It is a more informed approach.)}
\end{baihoc}
\end{truyen}

\section{Ngủ Gật Vì Đi Làm Thêm Nuôi Bản Thân}
\begin{truyen}{Ngủ Gật Vì Đi Làm Thêm Nuôi Bản Thân}{Falling Asleep Because of a Part-Time Job to Support Themselves}
\chuhoa{H}{ọc sinh ngủ gật trong lớp.} Giáo viên gõ bàn: ``Thức dậy. Em nghĩ lớp học là chỗ ngủ à?''

\textit{(A student falls asleep in class. The teacher taps the desk: ``Wake up. Do you think this classroom is a place to sleep?'')}

Học sinh mở mắt. Không xin lỗi. Nhìn thẳng: ``Thưa thầy, em đi làm từ sáu giờ tối đến hai giờ sáng để trả tiền nhà trọ.''

\textit{(The student opens their eyes. Does not apologise. Looks straight ahead: ``I work from six in the evening to two in the morning to pay rent.'')}

Cả lớp im. Giáo viên ngồi xuống. Một lúc sau nói: ``Cho thầy vài phút.''
\textit{(The class goes silent. The teacher sits down. After a moment: ``Give me a few minutes.'')}

\begin{baihoc}
Không phải mọi học sinh đều xuất phát từ cùng một điểm. Giáo viên nhìn thấy hành vi. Nhưng đằng sau hành vi đó có thể là một cuộc đời rất khác với những gì giáo viên tưởng. Hỏi một lần, trước khi phạt, có thể thay đổi tất cả.
\textit{(Not every student starts from the same position. Teachers see behaviour. But behind that behaviour there may be a life very different from what the teacher imagines. Asking once, before punishing, can change everything.)}
\end{baihoc}
\end{truyen}

\section{Bỏ Học Nhóm Vì Nhóm Không Có Chỗ Cho Mình}
\begin{truyen}{Bỏ Học Nhóm Vì Nhóm Không Có Chỗ Cho Mình}{Skipping Group Work Because the Group Has No Room for You}
\chuhoa{H}{ọc sinh không tham gia bài nhóm.} Giáo viên gọi lên hỏi. Học sinh nói: ``Tụi nó không phân công em việc gì. Em nhắn nhưng không ai trả lời.''
\textit{(A student does not participate in the group assignment. The teacher calls them over. The student says: ``They did not assign me any task. I messaged but nobody replied.'')}

\teacherpair{Thì em phải chủ động hơn chứ.}{``Then you need to be more proactive.''}

\studentpair{Em nhắn năm lần rồi thầy ơi. Tụi nó để em ngoài nhóm chat từ đầu.}{``I sent five messages, teacher. They excluded me from the group chat from the beginning.''}

Giáo viên dừng lại. Đây không còn là chuyện bài nhóm nữa.
\textit{(The teacher stops. This is no longer about the assignment.)}

\begin{baihoc}
Bài nhóm không chỉ kiểm tra kiến thức. Nó kiểm tra cả động học nhóm. Khi một học sinh liên tục bị loại khỏi nhóm, đó là vấn đề xã hội nghiêm trọng hơn nhiều so với điểm bài nhóm.
\textit{(Group work does not only test knowledge. It also tests group dynamics. When a student is consistently excluded, that is a social problem far more serious than the assignment grade.)}
\end{baihoc}
\end{truyen}

\section{Đến Trường Chỉ Để Không Ở Nhà}
\begin{truyen}{Đến Trường Chỉ Để Không Ở Nhà}{Coming to School Just to Not Be Home}
\chuhoa{C}{ó những học sinh không thích học} nhưng đến trường đều, đúng giờ, ngồi lặng lẽ và không gây chuyện. Không phải vì chúng ngoan. Mà vì nhà của chúng không phải nơi muốn ở.
\textit{(Some students do not enjoy studying but attend every day, on time, sit quietly, and cause no trouble. Not because they are well-behaved. But because their home is not a place they want to be.)}

Đối với những học sinh đó, trường học là nơi an toàn nhất trong ngày.
\textit{(For those students, school is the safest place in their day.)}

\begin{baihoc}
Khi giáo viên nghĩ về phòng học, họ thường nghĩ về nơi dạy và học. Nhưng với không ít học sinh, phòng học là nơi duy nhất trong ngày mà chúng cảm thấy ổn. Điều đó đặt lên vai người dạy một trách nhiệm lớn hơn chỉ là kiến thức.
\textit{(When teachers think about the classroom, they usually think about a place of teaching and learning. But for many students, the classroom is the only place in the day where they feel okay. That places on teachers a responsibility far greater than curriculum.)}
\end{baihoc}
\end{truyen}

\section{Trốn Trong Nhà Vệ Sinh Giữa Giờ Ra Chơi}
\begin{truyen}{Trốn Trong Nhà Vệ Sinh Giữa Giờ Ra Chơi}{Hiding in the Bathroom During Recess}
\chuhoa{M}{ột học sinh} không ra sân chơi. Luôn ở trong nhà vệ sinh suốt mười lăm phút giờ giải lao. Giám thị ghi vào sổ: ``Hay ở trong nhà vệ sinh. Nghi ngờ có vấn đề.''
\textit{(A student never goes out to the yard during recess. Always stays in the bathroom for the full fifteen minutes. The supervisor writes in the log: ``Often in the bathroom. Suspected problem.'')}

Không ai hỏi học sinh đó. Không ai biết rằng mỗi lần ra sân là nó bị một nhóm học sinh khác trêu chọc theo cách không để lại dấu vết.
\textit{(Nobody asks the student. Nobody knows that every time the student goes out, a group of other students harasses them in ways that leave no marks.)}

Giám thị điền vào ô: \textit{``Đã giải quyết''}. Mà không giải quyết gì cả.
\textit{(The supervisor checks the box: ``Resolved.'' Having resolved nothing.)}

\begin{baihoc}
Quan sát mà không hỏi chỉ tạo ra hồ sơ, không tạo ra sự giúp đỡ. Đứa học sinh đang trốn không cần ai ghi vào sổ. Nó cần ai đó dừng lại một chút và hỏi thật: ``Em có ổn không?''
\textit{(Observation without asking creates records, not help. The student in hiding does not need someone to write in a log. They need someone to stop and ask genuinely: ``Are you okay?'')}
\end{baihoc}
\end{truyen}

\begin{baihoc}
Trốn chạy không phải luôn là yếu đuối. Đôi khi, đó là cách duy nhất người ta còn có để tự bảo vệ mình khi không đủ can đảm hoặc không đủ ngôn từ để nói thật.

Người lớn có thể làm ít hơn rất nhiều thứ nếu chịu hỏi thêm một câu.
\textit{(Hiding is not always weakness. Sometimes it is the only form of self-protection left when a person lacks the courage or the words to speak truthfully. Adults could prevent a great deal simply by asking one more question.)}
\end{baihoc}

\begin{ghinhoanh}
\textbf{Short English to remember:}

Behaviour is communication. Learn to read it.

Sometimes the student doesn't need a solution. They need to be asked.
\end{ghinhoanh}

\begin{tuhocnhanh}
1. Em đã từng trốn khỏi một tình huống nào đó ở trường mà không có ai hỏi tại sao không? \textit{(Have you ever avoided a school situation that nobody ever asked you about?)}

2. Nếu em thấy bạn có hành vi lạ, em sẽ làm gì đầu tiên? \textit{(If you notice a friend behaving strangely, what is the first thing you would do?)}
\end{tuhocnhanh}

\ngancach
